\clearpage\SourceLeaf{316}

\begin{center}
[II.]

[THE PHYSICAL POSSIBILITY OF THE DATA IN POTHENOT'S PROBLEM.]

[1.]
\end{center}

\begin{center}
\textsc{Gauss} to \textsc{Schumacher}. Göttingen, 13 April 1836.
\end{center}

\indent \mbox{. . . . . . .} On this occasion I should like to mention another problem as well, namely the one known as \textsc{Pothenot}'s problem. If three points $A$, $B$, $C$ are given, together with the angles that $XA$, $XB$, $XC$ make with one another (including, of course, the \emph{sens}, i.e. whether the angles are reckoned from the left to the right [side] or conversely), then $X$ (generally speaking, i.e. except in the case where $A$, $B$, $C$, $X$ lie on one circle) is \textls[120]{determined}. If the data are based on an actual concrete case, it goes without saying that the problem, being real, is also possible. But if someone (to lead you astray) gives you only the numbers, it may turn out afterwards that the problem is physically impossible and that instead of one of the three directions you obtain the opposite direction (differing from the prescribed one by $180^0$).

The question is how \textls[120]{most simply} to express the condition of physical possibility by an equation \textls[120]{between the data}. Surprised that in a problem on which so much has been written, so essential a circumstance has, as far as I know, been entirely overlooked by everyone, I have often presented the solution in my lectures. Yet it probably cannot be altogether easy to find. At any rate, on one occasion I proposed it as the subject of a doctoral dissertation to a young man who was otherwise very capable (and to whom I had \textls[120]{not} presented the solution). Although he struggled with it for a long time, he could not find the solution. \mbox{. . . . . . .}
